\documentclass[10pt]{article}
\usepackage[a4paper,margin=0.88in]{geometry}
\usepackage[T1]{fontenc}
\usepackage{lmodern}
\usepackage{amsmath,amssymb,amsthm}
\usepackage{graphicx}
\usepackage{microtype}
\usepackage[colorlinks=true,linkcolor=blue,urlcolor=blue,citecolor=blue]{hyperref}
\newtheorem{theorem}{Theorem}
\newtheorem{lemma}{Lemma}
\newcommand{\R}{\mathbb R}
\newcommand{\Var}{\operatorname{Var}}
\title{Finite Exponential Moment Does Not Force\\
Finite-Time Ornstein--Uhlenbeck Poincar\'e Regularization}
\author{Candidate full counterexample packet\\
\small arXiv:2504.08658v2, Section 3.2.2}
\date{11 August 2026}

\begin{document}
\maketitle

\paragraph{Status.}
Candidate full negative answer, subject to expert review.  We construct a
normalized centered nonnegative $v\in H^1(\R,d\gamma)$ with a finite
exponential moment such that $|w(t)|^2d\gamma$ fails every Poincar\'e
inequality at every finite Ornstein--Uhlenbeck time $t$.

\section{Source question}
Let $d\gamma(x)=(2\pi)^{-1/2}e^{-x^2/2}\,dx$.  In Section 3.2.2,
Brigati--Dolbeault--Simonov \cite[p.~17]{BDS} consider the nonlinear form of
the Ornstein--Uhlenbeck evolution
\begin{equation}\label{eq:nonlinear}
 \partial_t w=w''+\frac{|w'|^2}{w}-xw',\qquad w(0)=v.
\end{equation}
They ask whether the assumption
\begin{equation}\label{eq:moment}
 v\in H^1(\R,d\gamma),\qquad
 \int_{\R}|v|^2e^{\theta|x|}\,d\gamma<\infty
 \quad\text{for some }\theta>0
\end{equation}
forces the existence of a finite time $T$ at which $|w(T)|^2d\gamma$
satisfies a Poincar\'e inequality.  Their surrounding discussion takes the
density to be normalized and centered; our example has both properties.

\begin{figure}[ht]
\centering
\includegraphics[width=.92\textwidth]{figures/open_question_crop.png}
\caption{The open question in the current source PDF, page 17.}
\end{figure}

\section{Construction and statement}
Put
\[
 a_k=2^k,\qquad Z=\sum_{k=1}^{\infty}e^{-a_k},\qquad
 p_k=Z^{-1}e^{-a_k}.
\]
Let $\nu_m=N(m,1)$ and define the symmetric probability measure
\begin{equation}\label{eq:mixture}
 \mu=\sum_{k=1}^{\infty}\frac{p_k}{2}
       \bigl(\nu_{a_k}+\nu_{-a_k}\bigr).
\end{equation}
Relative to $\gamma=N(0,1)$, the translate $\nu_m$ has density
\[
 g_m(x)=\exp(mx-m^2/2).
\]
Thus $\mu=f\gamma$, where
\begin{equation}\label{eq:density}
 f(x)=\sum_{k=1}^{\infty}p_k e^{-a_k^2/2}\cosh(a_kx),
 \qquad v=\sqrt f.
\end{equation}

\begin{theorem}\label{thm:main}
The function $v$ in \eqref{eq:density} is nonnegative, normalized, centered,
and belongs to $H^1(\R,d\gamma)$.  It satisfies \eqref{eq:moment} for every
$0<\theta<1$.  Nevertheless, for every finite $t\ge0$, the measure
$|w(t)|^2d\gamma$ generated by \eqref{eq:nonlinear} has no positive
Poincar\'e constant.  In particular the finite-time question of \cite{BDS}
has a negative answer.
\end{theorem}

\section{Regularity and exponential moment}
Normalization and centering follow at once from the probabilistic mixture
\eqref{eq:mixture}.  If $X$ has law $\mu$, then conditionally on the index and
sign it has the form $X=\pm a_k+G$, where $G\sim N(0,1)$.  Hence for
$0<\theta<1$,
\begin{align*}
 \mathbb E e^{\theta|X|}
 &\le \mathbb E e^{\theta|G|}\sum_{k=1}^{\infty}p_ke^{\theta a_k}\\
 &=\frac{\mathbb E e^{\theta|G|}}{Z}
   \sum_{k=1}^{\infty}e^{-(1-\theta)2^k}<\infty.
\end{align*}

It remains to verify the nontrivial $H^1$ requirement.  Index all signed
components by $i$, writing $f=\sum_iq_i g_{m_i}$.  Since
$g_{m_i}'=m_i g_{m_i}$, Cauchy--Schwarz gives pointwise
\begin{equation}\label{eq:fisher}
 \frac{|f'|^2}{f}
 =\frac{|\sum_iq_i g_{m_i}'|^2}{\sum_iq_i g_{m_i}}
 \le \sum_iq_i\frac{|g_{m_i}'|^2}{g_{m_i}}
 =\sum_iq_i m_i^2g_{m_i}.
\end{equation}
The series and its derivatives converge locally uniformly; alternatively,
\eqref{eq:fisher} follows first for finite mixtures and then by monotone
approximation.  Integrating and using $\int g_m\,d\gamma=1$ yields
\[
 4\int_{\R}|v'|^2\,d\gamma
 =\int_{\R}\frac{|f'|^2}{f}\,d\gamma
 \le\sum_{k=1}^{\infty}p_ka_k^2
 =\frac1Z\sum_{k=1}^{\infty}4^ke^{-2^k}<\infty.
\]
Together with $\int v^2d\gamma=1$, this proves $v\in H^1(\gamma)$.

\section{Exact evolution of the mixture}
Let $P_t$ be the Ornstein--Uhlenbeck semigroup and set $r=e^{-t}$.  Mehler's
formula gives the exact identity
\begin{equation}\label{eq:mehler}
 P_tg_m(x)=\exp(rmx-r^2m^2/2)=g_{rm}(x).
\end{equation}
Consequently $\mu_t=(P_tf)\gamma$ is the symmetric mixture
\begin{equation}\label{eq:evolved}
 \mu_t=\sum_{k=1}^{\infty}\frac{p_k}{2}
       \bigl(N(ra_k,1)+N(-ra_k,1)\bigr).
\end{equation}
If $w(t)=\sqrt{P_tf}$, then
$\partial_t(w^2)= (\partial_x^2-x\partial_x)(w^2)$, and division by $2w$
gives \eqref{eq:nonlinear}.  Thus \eqref{eq:evolved} is exactly the measure
appearing in the source question.

\section{Vanishing Poincar\'e Rayleigh quotients}
Fix an arbitrary finite $t$, so $r>0$, and write $b_j=ra_j$.  For $k\ge2$ set
\[
 L_k=\frac{b_{k-1}+b_k}{2},\qquad
 U_k=\frac{b_k+b_{k+1}}2.
\]
Choose a smooth cutoff $h_k$ with values in $[0,1]$ which is one on
$[L_k+1,U_k-1]$, zero outside $[L_k-1,U_k+1]$, and satisfies
$|h_k'|\le1$.  Such cutoffs exist for all sufficiently large $k$.

Under the component $N(b_k,1)$, the plateau endpoints are at signed
distances
\[
 -\frac{r(a_k-a_{k-1})}{2}+1=-r2^{k-2}+1,\qquad
 \frac{r(a_{k+1}-a_k)}2-1=r2^{k-1}-1
\]
from its mean.  Therefore this component assigns the plateau probability at
least $1/2$ for all large $k$.  Hence
\[
 \mu_t\{h_k=1\}\ge\frac{p_k}{4}.
\]
On the other hand, at least half of every negative-mean Gaussian lies to the
left of $L_k-1$ for large $k$.  The total negative-component weight is $1/2$,
so $\mu_t\{h_k=0\}\ge1/4$.  Using
$\Var_{\mu_t}h=\frac12\iint|h(x)-h(y)|^2d\mu_t(x)d\mu_t(y)$, we obtain
\begin{equation}\label{eq:variance}
 \Var_{\mu_t}(h_k)\ge\frac{p_k}{16}.
\end{equation}

The derivative is supported in the two intervals of radius one centered at
$L_k$ and $U_k$.  Every component mean $\pm b_j$ is at distance at least
$r2^{k-2}$ from their union's nearer midpoint.  A standard normal assigns an
interval of length two at midpoint-distance $d>1$ mass at most
$2(2\pi)^{-1/2}e^{-(d-1)^2/2}$.  Summing all mixture weights therefore gives
\begin{equation}\label{eq:energy}
 \int|h_k'|^2d\mu_t
 \le \frac4{\sqrt{2\pi}}
       \exp\!\left[-\frac12\bigl(r2^{k-2}-1\bigr)^2\right].
\end{equation}
Combining \eqref{eq:variance}, \eqref{eq:energy}, and
$p_k=Z^{-1}e^{-2^k}$,
\begin{equation}\label{eq:rayleigh}
 \frac{\int|h_k'|^2d\mu_t}{\Var_{\mu_t}(h_k)}
 \le \frac{64Z}{\sqrt{2\pi}}
 \exp\!\left[2^k-\frac12\bigl(r2^{k-2}-1\bigr)^2\right]
 \longrightarrow0.
\end{equation}
The cutoff can be chosen directly smooth as above, or obtained by smoothing
a piecewise linear cutoff without changing the estimates.  Thus no constant
$C_P>0$ can make
\[
 \int|\phi'|^2d\mu_t\ge C_P\Var_{\mu_t}(\phi)
 \quad\text{for all }\phi\in C_c^\infty(\R).
\]
Because the finite time $t$ was arbitrary, this proves Theorem
\ref{thm:main}.

\section{Scope and verification}
The counterexample is one-dimensional, which suffices to refute the
dimension-free question.  It exploits infinitely many increasingly remote
bottlenecks.  Its finite exponential moment is strictly weaker than the
source's Gaussian two-point moment condition (22), which this mixture does
not satisfy, so it does not conflict with the cited positive theorem.

The formula in \eqref{eq:rayleigh} was checked numerically on a logarithmic
scale for several finite times.  A bounded search through 11 August 2026 of
the run registry, the current arXiv source and published DOI, exact phrases,
and Gaussian-mixture/Poincar\'e keywords found no later resolution or this
construction.  This is not a claim of priority.

\paragraph{Human-review recommendation.}
Verify the transition-strip bound in \eqref{eq:energy}.  Also check the
identity $|w(t)|^2d\gamma=\mu_t$.  The obstruction is then immediate.

\begin{thebibliography}{9}
\bibitem{BDS}
G. Brigati, J. Dolbeault, and N. Simonov,
\emph{Logarithmic Sobolev inequalities: a review on stability and instability
results; logarithmic uncertainty principle}, La Matematica 5 (2026), article
5; \href{https://arxiv.org/abs/2504.08658}{arXiv:2504.08658v2}.
\end{thebibliography}
\end{document}
